\documentclass[a4paper,12pt]{article}
\usepackage[brazil]{babel}
\usepackage[utf8]{inputenc}
\usepackage[T1]{fontenc}
\usepackage{geometry}
\usepackage{setspace}
\usepackage{titlesec}
\usepackage{hyperref}
\usepackage{graphicx}
\usepackage{caption}
\usepackage{subcaption}
\usepackage{fancyhdr}
\usepackage{xcolor}
\usepackage{amsmath,amssymb,physics}
\usepackage{graphicx}
\usepackage{hyperref}
\usepackage{bm}
\usepackage{empheq}
\usepackage[most]{tcolorbox}


\input{commands.tex}

\hypersetup{
    colorlinks=true,% make the links colored
    linkcolor=blue
}

\usepackage{setspace}
\addbibresource{bibliography.bib}

\newcommand{\printingbibliography}{%

    \pagestyle{myheadings}
    \markright{}
    \sloppy
    \printbibliography[heading=bibintoc, % add to table of contents
                   title=Refer\^encias % Chapter name
                  ]
    \fussy%
}
\PassOptionsToPackage{table}{xcolor}
\renewcommand{\baselinestretch}{1.5}

\fancyhf{}
\fancyhead[L]{Andr\'e Vieira da Silva}
\fancyhead[R]{Algoritmos de Otimização em Computação Quântica \thepage}
\renewcommand{\headrulewidth}{0.8pt}
\pagestyle{fancy}

\usepackage{cleveref}
\newtcbtheorem{Theorem}{Theorem}{
  enhanced,
  sharp corners,
  attach boxed title to top left={
    yshifttext=-1mm
  },
  colback=white,
  colframe=blue!75!black,
  fonttitle=\bfseries,
  boxed title style={
    sharp corners,
    size=small,
    colback=blue!75!black,
    colframe=blue!75!black,
  } 
}{thm}
\newtcbtheorem[no counter]{Proof}{}{
  enhanced,
  sharp corners,
  attach boxed title to top left={
    yshifttext=-1mm
  },
  colback=white,
  colframe=blue!25,
  fonttitle=\bfseries,
  coltitle=black,
  boxed title style={
    sharp corners,
    size=small,
    colback=blue!25,
    colframe=blue!25,
  } 
}{prf}

\geometry{a4paper,top=30mm,bottom=20mm,left=30mm,right=20mm}

\titleformat*{\section}{\bfseries\large}
\titleformat*{\subsection}{\bfseries\normalsize}

\title{Algoritmos Variacionais e de Otimização em Computação Quântica \\ Andr\'e V. Silva}
\author{\textcolor{blue}{DESAFIO PROCESSO SELETIVO ITAÚ} \\ \textcolor{blue}{BOLSISTA TECH 3 MESTRE PROJ1220120573}}
\date{\today}

\begin{document}

\begin{Proof}{Desafio}{}
"Explicar em detalhes a construção e o funcionamento do algoritmo de VQE e QAOA, QWOA, BF-DCQO.
Explicando todas as etapas e a influência dos Gates na construção desses algoritmos. 
Discutir os benefícios e desafios no contexto desses algoritmos."
\end{Proof}

\maketitle
\thispagestyle{fancy}

\tableofcontents

\newpage

\section{Introdução}

A \colorbox{yellow!30}{computação quântica}\cite{drinko2024variational} tem se desenvolvido rapidamente nas últimas décadas 
como uma nova forma de \colorbox{blue!20}{processamento de informação} baseada nos princípios 
fundamentais da \colorbox{blue!20}{mecânica quântica}\cite{cohen2019quantum, sakurai2020modern}. Diferentemente dos \colorbox{green!20}{computadores clássicos}, 
que utilizam bits que assumem valores \colorbox{green!20}{binários ($0$ ou $1$)}, os \colorbox{red!20}{computadores 
quânticos} utilizam \colorbox{red!20}{\textit{qubits}}(Quantum Bits), que podem existir em \colorbox{yellow!20}{superposição de estados}.

Um qubit pode ser descrito pelo estado

\begin{equation}
|\psi\rangle = \alpha |0\rangle + \beta |1\rangle = 
\alpha
\begin{pmatrix}
1 \\
0
\end{pmatrix}
+
\beta
\begin{pmatrix}
0 \\
1
\end{pmatrix}
=
\begin{pmatrix}
\alpha \\
\beta
\end{pmatrix}
\end{equation}

onde $\alpha$ e $\beta$ são amplitudes complexas que satisfazem a \colorbox{yellow!20}{condição de normalização}

\begin{equation}
\boxed{
|\alpha|^2 + |\beta|^2 = 1
}
\end{equation}

Além da superposição, sistemas quânticos também podem apresentar \colorbox{red!20}{\textit{emaranhamento}}
e \colorbox{red!20}{\textit{interferência}}, propriedades que permitem novas formas de processamento 
de informação.

Entretanto, os \colorbox{blue!20}{computadores quânticos atuais pertencem à chamada era \textit{NISQ}} 
(\textit{Noisy Intermediate-Scale Quantum}), caracterizada por:

\begin{itemize}
\item número limitado de qubits
\item presença significativa de ruído
\item profundidade limitada de circuitos quânticos
\end{itemize}

Devido a essas \colorbox{red!20}{limitações}, foram desenvolvidos \colorbox{blue!20}{algoritmos híbridos quântico-clássicos}
que combinam \colorbox{blue!20}{circuitos quânticos parametrizados} com \colorbox{green!20}{otimização clássica}. 
Entre esses algoritmos destacam-se:

\begin{itemize}
\item \colorbox{yellow!30}{Variational Quantum Eigensolver (VQE)}
\item \colorbox{yellow!30}{Quantum Approximate Optimization Algorithm (QAOA)}
\item \colorbox{yellow!30}{Quantum Whale Optimization Algorithm (QWOA)}
\item \colorbox{yellow!30}{Bias-Field Digitized Counterdiabatic Quantum Optimization (BF-DCQO)}
\end{itemize}

Esses \colorbox{green!20}{algoritmos exploram circuitos quânticos parametrizados} para resolver problemas 
de \colorbox{green!20}{otimização, simulação quântica e busca por estados fundamentais.}

\section{Circuitos Quânticos Parametrizados}

Os algoritmos variacionais utilizam circuitos chamados \colorbox{blue!20}{\textit{Parameterized 
Quantum Circuits}} (PQC). Esses circuitos possuem \colorbox{blue!20}{portas quânticas dependentes de 
parâmetros contínuos}, geralmente \colorbox{blue!20}{ângulos de rotação}.

Um exemplo típico de porta parametrizada é

\begin{equation}
\boxed{
R_y(\theta) = e^{-i\theta Y/2}
}
\end{equation}

onde $Y$ é a matriz de Pauli:

\begin{equation}
\sigma_x =
\begin{pmatrix}0 & 1 \\ 1 & 0\end{pmatrix},
\quad
\sigma_y =
\begin{pmatrix}0 & -i \\ i & 0\end{pmatrix},
\quad
\sigma_z =
\begin{pmatrix}1 & 0 \\ 0 & -1\end{pmatrix}.
\end{equation}


O estado produzido por um \colorbox{blue!20}{circuito parametrizado} pode ser escrito como

\begin{equation}
|\psi(\textcolor{red}{\theta})\rangle = \textcolor{blue}{U}(\textcolor{red}{\theta})|0\rangle^{\otimes n} = \textcolor{blue}{U}(\textcolor{red}{\theta})|\underbrace{000...00}_{n}\rangle
\end{equation}

onde $\textcolor{blue}{U}(\textcolor{red}{\theta})$ representa o circuito quântico dependente do vetor de 
parâmetros $\textcolor{red}{\theta}$.

\colorbox{green!30}{O objetivo é encontrar valores ótimos desses parâmetros minimizando uma função}\newline 
\colorbox{green!30}{custo}.

\section{Variational Quantum Eigensolver (VQE)}

\begin{tcolorbox}[colback=red!10,colframe=red!70!black,title=VQE]
Algoritmo usado principalmente para 
encontrar o estado de menor energia de um sistema quântico.
\end{tcolorbox}

O \colorbox{green!30}{algoritmo VQE} foi desenvolvido para \colorbox{green!30}{calcular autovalores de Hamiltonianos}, 
especialmente a \colorbox{green!30}{energia fundamental de sistemas quânticos}.

Considere um \colorbox{blue!30}{Hamiltoniano $H$ que descreve um sistema físico}. O problema 
consiste em \colorbox{red!30}{encontrar o menor autovalor da equação}

\begin{equation}
H|\psi\rangle = E|\psi\rangle,
\end{equation}

onde:

\begin{itemize}
\item $H$ → Hamiltoniano do sistema
\item $E$ → energia
\item $|\psi \rangle$ → estado quântico
\end{itemize}

O \colorbox{green!30}{VQE} baseia-se no \colorbox{green!30}{princípio variacional da mecânica quântica}, que afirma que

\begin{equation}
E(\theta) = \langle \psi(\theta)|H|\psi(\theta)\rangle \geq E_0
\end{equation}

onde \colorbox{green!30}{$E_0$ é a energia do estado fundamental}.

\subsection{Etapas do VQE}

O funcionamento do algoritmo envolve as seguintes etapas:

\begin{enumerate}

\item Escolher um ansatz parametrizado para o estado quântico.

\item Preparar o estado

\begin{equation}
|\psi(\theta)\rangle
\end{equation}

utilizando portas de rotação e portas de emaranhamento.

\item Medir o valor esperado do Hamiltoniano

\begin{equation}
E(\theta)=\langle \psi(\theta)|H|\psi(\theta)\rangle
\end{equation}

\item Utilizar um otimizador clássico para atualizar os parâmetros $\theta$.

\item Repetir o processo até convergência.

\end{enumerate}

\subsection{Decomposição do Hamiltoniano}

Em implementações práticas, o Hamiltoniano é decomposto como

\begin{equation}
H = \sum_i c_i P_i
\end{equation}

onde $P_i$ são produtos de matrizes de Pauli.

Cada termo é medido separadamente no computador quântico.

\subsection{Por que o VQE é importante?}

Ele é um dos algoritmos mais promissores para computadores NISQ.

Aplicações:
\begin{itemize}
\item química quântica
\item estrutura eletrônica molecular
\item otimização combinatória
\item física do estado sólido.
\end{itemize}

\section{Quantum Approximate Optimization Algorithm (QAOA)}

\begin{tcolorbox}[colback=red!10,colframe=red!70!black,title=QAOA]
O QAOA foi criado para resolver problemas de otimização combinatória usando 
computadores qu\^anticos.
\end{tcolorbox}

Ele utiliza dois Hamiltonianos principais:

\begin{itemize}
\item Hamiltoniano do problema $H_C$
\item Hamiltoniano mixer $H_M$
\end{itemize}

O estado inicial geralmente é

\begin{equation}
|\psi_0\rangle = |+\rangle^{\otimes n}
\end{equation}

obtido aplicando portas de Hadamard.

A evolução do sistema é dada por

\begin{equation}
|\psi(\gamma,\beta)\rangle =
\prod_{k=1}^{p} e^{-i\beta_k H_M} e^{-i\gamma_k H_C} |\psi_0\rangle
\end{equation}

Os parâmetros $\gamma_k$ e $\beta_k$ são ajustados por um otimizador clássico.

O valor esperado do Hamiltoniano de custo é então minimizado.



\section{Quantum Whale Optimization Algorithm (QWOA)}

\begin{tcolorbox}[colback=red!10,colframe=red!70!black,title=QWOA]
O QWOA é uma adaptação quântica do algoritmo clássico Whale Optimization Algorithm, 
inspirado no comportamento de caça das baleias jubarte. 
\end{tcolorbox}

Nesse algoritmo, a população de soluções é representada por estados quânticos 
em superposição

\begin{equation}
|\psi\rangle = \sum_x \alpha_x |x\rangle
\end{equation}

A evolução do algoritmo envolve operadores quânticos que simulam:

\begin{itemize}
\item aproximação da presa
\item exploração global
\item movimento em espiral
\end{itemize}

Essas operações são implementadas por meio de portas quânticas parametrizadas e 
operadores controlados que modificam as amplitudes dos estados.

O processo iterativo altera as amplitudes de probabilidade de forma a aumentar a 
probabilidade de observar soluções ótimas.

\section{Bias-Field Digitized Counterdiabatic Quantum Optimization (BF-DCQO)}

\begin{tcolorbox}[colback=red!10,colframe=black!70!black,title=BF-DCQO]
O BF-DCQO é um algoritmo baseado em evolução adiabática digitalizada.
\textbf{Ideia simples}: Ele tenta guiar o sistema quântico até o estado de menor 
energia evitando erros durante a evolução.
\end{tcolorbox}

O BF-DCQO é um algoritmo inspirado em métodos de otimização baseados em 
\textit{quantum annealing}.

Nesse tipo de algoritmo, o sistema evolui sob um \colorbox{yellow!20}{Hamiltoniano dependente do tempo}

\begin{equation}
H(t)=A(t)H_i + B(t)H_p
\end{equation}

onde:

\begin{itemize}
\item $H_i$ é o Hamiltoniano inicial
\item $H_p$ é o Hamiltoniano do problema
\end{itemize}

Durante a evolução adiabática podem ocorrer transições para estados excitados. 
Para evitar esse problema, introduz-se um termo contradiabático

\begin{equation}
H_{CD}(t)
\end{equation}

O Hamiltoniano total torna-se

\begin{equation}
H(t)=A(t)H_i + B(t)H_p + C(t)H_{CD}
\end{equation}

Esse termo reduz transições não adiabáticas e melhora a probabilidade de encontrar 
o estado fundamental.

Em implementações digitais, a evolução contínua é discretizada utilizando a 
decomposição de Trotter

\begin{equation}
U = \prod_k e^{-iH_k \Delta t}
\end{equation}

permitindo a implementação por meio de circuitos quânticos compostos por portas elementares.

\section{Benefícios e Desafios}

Esses algoritmos apresentam diversas vantagens no contexto da computação quântica atual.

Entre os principais benefícios destacam-se:

\begin{itemize}
\item compatibilidade com dispositivos NISQ
\item utilização de circuitos relativamente rasos
\item aplicações em química quântica e otimização
\end{itemize}

Entretanto, ainda existem desafios importantes, tais como:

\begin{itemize}
\item ruído em hardware quântico
\item grande número de medições necessárias
\item presença de barren plateaus na otimização
\item dificuldade na escolha de ansatz eficientes
\end{itemize}

\section{Conclusão}

\begin{tcolorbox}[colback=blue!10,colframe=blue!70!black,title=Algoritmos]
\begin{itemize}

\item São algoritmos de otimização
    \begin{itemize}
        \item Todos procuram minimizar ou maximizar uma função custo
    \end{itemize}

\item Usam circuitos quânticos parametrizados
    \begin{itemize}
        \item Eles usam portas quânticas com parâmetros ajustáveis
    \end{itemize}

\item São algoritmos híbridos (clássico + quântico)

\item Funcionam bem em computadores quânticos atuais (NISQ)
    \begin{itemize}
        \item poucos qubits
        \item muito ruído
        \item circuitos curtos
    \end{itemize}

\end{itemize}
\end{tcolorbox}

Os métodos \colorbox{green!20}{VQE, QAOA, QWOA e BF-DCQO utilizam circuitos parametrizados} e otimização 
híbrida para resolver problemas complexos em física, química e ciência da computação.

Apesar das limitações impostas pelo ruído e pela escalabilidade dos dispositivos 
quânticos atuais, esses algoritmos continuam sendo intensamente estudados e aprimorados, 
com o objetivo de alcançar vantagens computacionais significativas em aplicações práticas.


%%%%%%%% Bibliography 
% Os comandos para incluir as referências bibliográficas
\printingbibliography

\end{document}
